\emph{Aan alle idioten, die zijn of nog zullen komen, TABÉ.}

\emph{De volledige Volkvergadering heeft unaniem en met één stem
besloten:}

\begin{center}\rule{0.5\linewidth}{0.5pt}\end{center}

\hypertarget{art.-1}{%
\paragraph{Art. 1}\label{art.-1}}

Artikel 36 van de Grondwet zal vanaf nu lezen als volgt:

\emph{Er is één Hof voorgezeten door de Linker. Zij behandelt alle zaken
in eerste aanleg met uitzondering deze vastgelegd bij wet.}

\emph{Er is één niet-permanent Tribunaal. Zij behandelt alle zaken
waarvoor de wet dit voorschrijft en alle voorzieningen in cassatie.}

\emph{Het Tribunaal is voorgezeten door de Linker, of een ander lid van
de Volksvergadering zoals vastgelegd bij wet, en twee bijzitters
aangeduid zoals vastgelegd bij wet.}

\emph{Het Tribunaal wordt voor elke zaak of voorziening opnieuw
samengesteld.}

\emph{Bij beroep kan niemand die een uitspraak heeft gedaan in eerste
aanleg de zaak opnieuw berechten, zij het als lid van een Tribunaal of
alleen.}

\hypertarget{art.-2}{%
\paragraph{Art. 2}\label{art.-2}}

Volgend artikel wordt toegevoegd aan de Grondwet:

\hypertarget{art.-36bis}{%
\subparagraph{\texorpdfstring{\emph{Art.
36bis}}{Art. 36bis}}\label{art.-36bis}}

\emph{Het Hof vonnist via vonnissen.}

\emph{Het Tribunaal vonnist via arresten.}

\hypertarget{art.-3}{%
\paragraph{Art. 3}\label{art.-3}}

Artikel 39 van de Grondwet zal vanaf nu lezen als volgt:

\emph{Men kan in beroep gaan tegen een vonnis.}

\emph{Men kan in cassatie gaan tegen een arrest.} \emph{Een nieuw
Tribunaal wordt hiervoor samengesteld zoals bedoeld in artikel 36.}
\emph{Dit Tribunaal kan enkel het betrokken arrest verbreken indien dit
in strijd is met de wet.} \emph{Dit Tribunaal spreekt zich niet uit over
de inhoud van zaken.}

\begin{center}\rule{0.5\linewidth}{0.5pt}\end{center}

\emph{De Kontsuls van het Jaar des Herens Tweeduizendzesentwintig}

\emph{Op den \texttt{\textless{}Datum\textgreater{}}}
